\chapter{Автоматизированная инфраструктура валидаторов}\label{ch:ch4}

\section{Требования и архитектурные принципы}\label{sec:ch4/sec1}

Формулируются нефункциональные требования к автоматизации: надежность, безопасность,
устойчивость к отказам, наблюдаемость и управляемость. Описываются архитектурные
принципы, обеспечивающие масштабирование и повторное использование валидаторов.

\section{Схема подбора валидаторов на основе производительности}\label{sec:ch4/sec2}

Описываются критерии отбора валидаторов, набор метрик производительности и
принципы динамического обновления состава сети. Отдельно выделяются требования к
прозрачности и воспроизводимости процедуры отбора.

\section{Методология экспериментов и моделирования}\label{sec:ch4/sec3}

Формируется план экспериментов: тестовая среда, сценарии нагрузки, методика
сбора метрик и критерии оценки качества автоматизации.
